
The validation confirms that published technical reports from three major LLM labs systematically contain category-level error rate data with granularity and completeness meeting predefined thresholds. This finding establishes that such data are available for potential downstream applications, though those applications were not tested in this study.

\textbf{Interpretation.} The consistency across three independent organizations suggests that detailed benchmark reporting has become standard practice, possibly driven by competitive pressure to demonstrate fine-grained performance. The uniformity in category structures (identical counts across families for each benchmark) indicates labs likely use common evaluation tools or adhere to benchmark-defined categories.

The availability of data across both baseline and current model generations enables potential longitudinal analysis, though such analysis was not performed in this study. One could examine which categories show improvement and which remain challenging, though interpreting such patterns would require additional validation.

\textbf{Limitations.} Several limitations constrain the scope of these findings:

\begin{enumerate}
\item \textit{Foundation-only validation.} This study validates that category-level data exist but does not demonstrate their sufficiency for specific applications. Potential uses such as error taxonomy generation, pattern clustering, or weak supervision approaches remain untested. Whether the available data support such applications is an open question requiring dedicated validation.
\end{enumerate}

\begin{enumerate}
\item \textit{Manual extraction method.} Using curated manual extraction rather than automated parsing affects scalability. Extending to additional model families or tracking data over longer time periods would benefit from automated extraction, which would likely achieve lower completeness (85-95\% based on pilot attempts).
\end{enumerate}

\begin{enumerate}
\item \textit{Published data accuracy assumed.} We treat published benchmark results as accurate without independent validation. Labs could report inaccurate results due to implementation errors or sampling issues. However, re-evaluation requires substantial compute resources and API access. The consistency across three independent labs with different incentives provides some confidence, though formal validation would strengthen this assumption.
\end{enumerate}

\begin{enumerate}
\item \textit{Limited scope.} Coverage is restricted to three frontier labs and two established benchmarks. Smaller labs, open-source community models, or newer benchmarks may show different reporting patterns. The findings may not generalize beyond the cases examined.
\end{enumerate}

\textbf{Implications.} If researchers can access category-level data without re-evaluation, this reduces barriers to error analysis. Researchers without API budgets or proprietary access could conduct analyses using publicly available data. However, realizing this potential requires validating that specific analysis methods (e.g., pattern extraction, clustering, expert validation) work with category-level granularity. Such validation is outside the scope of this study.

\textbf{Dependence on industry practices.} The availability of published data depends on labs' continued willingness to disclose category-level results. Changes in disclosure practices, competitive dynamics, or regulatory environments could reduce data availability. Building community-run evaluation infrastructure could reduce dependence on industry disclosures.
